\documentclass[11pt,a4paper]{article}

% ---------------------------------------------------------------------------
% Packages
% ---------------------------------------------------------------------------
\usepackage[utf8]{inputenc}
\usepackage[T1]{fontenc}
\usepackage[margin=1in]{geometry}
\usepackage{booktabs}
\usepackage{longtable}
\usepackage{array}
\usepackage{enumitem}
\usepackage{xcolor}
\usepackage{titlesec}
\usepackage{amsmath}
\usepackage{amssymb}
\usepackage[hidelinks]{hyperref}

\titleformat{\section}{\large\bfseries\color{black!85}}{\thesection}{0.6em}{}
\titleformat{\subsection}{\normalsize\bfseries}{\thesubsection}{0.6em}{}
\setlist{itemsep=2pt,topsep=3pt}
\renewcommand{\arraystretch}{1.25}

\newcommand{\dsname}[1]{\textbf{#1}}

% ---------------------------------------------------------------------------
\title{\vspace{-1.5em}\textbf{Generative-AI--Based Automated Analysis and\\
Risk Scoring of Fraudulent APKs}\\[0.4em]
\large Datasets and Solution Approach}
\author{}
\date{}

\begin{document}
\maketitle
\vspace{-3em}

% ===========================================================================
\section{Introduction}
% ===========================================================================
Android remains the dominant mobile platform and, consequently, the primary target of
mobile-banking fraud. Malicious APKs increasingly impersonate legitimate banking,
payment, and utility applications to steal one-time passwords (OTPs), harvest
credentials, abuse accessibility services, and exfiltrate sensitive data. Manual triage
of such samples is slow and does not scale with the volume of submissions a fraud or
Security Operations Centre (SOC) team receives.

This document describes the \emph{datasets we will use} to train and validate the
machine-learning component of our platform, and the \emph{solution approach} we will
follow to deliver automated, explainable risk assessment of suspicious APKs. Our system
combines transparent rule-based scoring, supervised machine learning, controlled dynamic
behaviour analysis, and a generative-AI reasoning layer that produces human-readable
investigation reports.

% ===========================================================================
\section{Datasets}
% ===========================================================================
The machine-learning model will be trained, validated, and stress-tested on a
combination of established Android-malware corpora and curated benign application sets.
We deliberately combine \emph{older, well-labelled} datasets (for breadth of malware
families and reproducibility) with \emph{recent} datasets and live sampling (to mitigate
concept drift, since malware behaviour evolves over time).

\subsection{Primary training and benchmark datasets}

\begin{itemize}
  \item \dsname{Drebin.}
  A widely cited Android-malware benchmark containing approximately \textbf{5{,}560
  malicious} samples spanning \textbf{179 malware families}, alongside a large set of
  benign applications. Crucially, Drebin provides \emph{pre-extracted static feature
  sets} per sample---requested and used permissions, API calls, intent actions, hardware
  components, application components, and embedded network indicators---together with a
  family-label mapping. These ready-made feature vectors allow us to train and benchmark
  classifiers without reprocessing raw APKs, making it our primary baseline dataset.

  \item \dsname{CICMalDroid 2020.}
  A modern dataset of roughly \textbf{17{,}000} Android samples organised into five
  classes: \emph{Adware}, \emph{Banking malware}, \emph{SMS malware}, \emph{Riskware},
  and \emph{Benign}. It includes both static attributes and dynamically observed
  behaviour (system calls, binder activity, composite behaviours). Its banking- and
  SMS-malware categories align directly with our fraud-triage objective, and its recency
  helps the model generalise beyond the older Drebin distribution.

  \item \dsname{CIC-AndMal2017.}
  A dataset emphasising \emph{dynamic and network-traffic} features, comprising several
  hundred malware samples across major families (ransomware, scareware, SMS malware,
  adware) and several thousand benign applications. We will use it to enrich the feature
  space with runtime and network-flow characteristics that static analysis cannot
  capture.

  \item \dsname{AndroZoo.}
  A continuously growing repository of \textbf{millions} of real-world APKs collected
  from Google Play and numerous third-party markets, each annotated with multi-engine
  anti-virus detection counts and collection metadata. We will use AndroZoo as a source
  of \emph{fresh, diverse} benign and malicious samples drawn from recent time windows,
  specifically to build an out-of-distribution evaluation set for measuring concept
  drift and real-world false-positive rates.

  \item \dsname{Android Malware Genome Project (MalGenome).}
  A historical reference set of \textbf{1{,}260} malware samples across \textbf{49
  families}. Although dated, it provides a clean, well-documented baseline of classic
  malware behaviours that is useful for sanity-checking detectors and for family-level
  analysis.

  \item \dsname{KronoDroid.}
  A hybrid dataset that records \emph{timestamped} samples with both static and dynamic
  features captured on real devices and emulators. Its temporal labelling makes it
  particularly valuable for time-aware evaluation, i.e.\ training on earlier samples and
  testing on later ones to simulate realistic deployment and quantify drift.
\end{itemize}

\subsection{Benign application sources}
A robust classifier requires a representative, modern benign distribution to keep
false-positive rates low. In addition to the benign partitions of the datasets above, we
will curate benign applications from:
\begin{itemize}
  \item \dsname{F-Droid} --- a catalogue of free and open-source Android applications,
  useful because the sources are auditable and the apps are demonstrably non-malicious.
  \item \dsname{Google Play (top and long-tail apps via AndroZoo)} --- to capture the
  permission and API profiles of mainstream commercial applications, including
  advertisement- and analytics-heavy apps that are benign but superficially
  ``suspicious.''
\end{itemize}

\subsection{Threat-intelligence corroboration source}
For ground-truth verification and IOC enrichment (not as a training feature), we will
consult \dsname{VirusTotal}, which aggregates verdicts from many anti-malware engines.
This supports label validation and provides an independent second opinion during
evaluation.

\subsection{Summary of datasets}
\begin{longtable}{p{3.0cm} p{2.0cm} p{3.0cm} p{4.5cm}}
\toprule
\textbf{Dataset} & \textbf{Approx.\ scale} & \textbf{Feature type} & \textbf{Role in our work}\\
\midrule
\endhead
Drebin & 5.5k mal.\ + benign & Static (perms, API, intents, IOCs) & Primary baseline; pre-extracted features for fast training\\
CICMalDroid 2020 & $\sim$17k & Static + dynamic & Modern multi-class (incl.\ banking/SMS) training\\
CIC-AndMal2017 & $\sim$5.5k & Dynamic + network & Runtime/network feature enrichment\\
AndroZoo & Millions & Raw APKs + AV labels & Fresh sampling; drift / false-positive evaluation\\
MalGenome & 1.26k & Static & Historical baseline; family sanity checks\\
KronoDroid & Tens of thousands & Hybrid, timestamped & Time-aware (drift) evaluation\\
F-Droid / Play & Curated & Raw APKs & Modern benign distribution\\
VirusTotal & On-demand & Multi-engine verdicts & Label validation / IOC corroboration\\
\bottomrule
\caption{Datasets to be used for model training, benchmarking, and evaluation.}
\end{longtable}

\subsection{Data preparation and partitioning}
\begin{itemize}
  \item \textbf{Normalisation and labelling.} Sample identities will be tracked by
  cryptographic hash; labels will be reconciled across datasets and corroborated with
  multi-engine verdicts to reduce label noise.
  \item \textbf{Class balance.} Because benign samples vastly outnumber malicious ones,
  we will apply class weighting and/or balanced resampling to avoid a model biased toward
  the majority class.
  \item \textbf{Splitting strategy.} We will use stratified train/validation/test splits
  for in-distribution evaluation and, additionally, a \emph{temporal split} (train on
  older samples, test on newer ones) to obtain an honest estimate of real-world
  performance under concept drift.
  \item \textbf{Leakage control.} Family-aware and time-aware splitting will be used to
  prevent near-duplicate samples from the same campaign appearing in both training and
  test sets, which would otherwise inflate reported accuracy.
\end{itemize}

% ===========================================================================
\section{Solution Approach}
% ===========================================================================
Our approach treats APK risk assessment as a layered pipeline. Each layer contributes
independent, explainable evidence, and the layers are fused into a single, transparent
risk score accompanied by a natural-language report. Generative AI is used to
\emph{reason over and communicate} the extracted evidence, never as the sole detector.

\subsection{Stage 1 --- Intake and identity}
Each submitted APK will be hashed (e.g.\ SHA-256/MD5) for de-duplication and blacklisting,
and basic metadata (package name, version, target/minimum SDK, file size) will be
recorded. Hash-based identity enables consistent labelling and correlation with
threat-intelligence feeds.

\subsection{Stage 2 --- Static analysis}
Static analysis inspects the application without executing it and forms the backbone of
fast, safe triage. We will extract:
\begin{itemize}
  \item \textbf{Manifest and permissions:} declared components, intent filters, and the
  set of requested permissions, with emphasis on fraud-relevant permissions
  (SMS, contacts, accessibility, overlay, install-packages, phone state).
  \item \textbf{Suspicious API usage:} references to APIs associated with OTP theft,
  accessibility abuse, dynamic code loading, command execution, and network
  communication, weighted by how discriminative each is in benign versus malicious apps.
  \item \textbf{Indicators of compromise (IOCs):} URLs, IP addresses, domains, embedded
  tokens, and likely secrets, including indicators recovered by decoding obfuscated
  (e.g.\ Base64/hex) strings and by scanning bundled native libraries.
  \item \textbf{Obfuscation and packing signals:} identifier-renaming ratios, multiple
  DEX files, high-entropy blobs, dynamic loaders, and packer/compiler fingerprints.
  \item \textbf{Signing certificate and attack surface:} certificate red-flags
  (debug/placeholder identities, implausible validity) and exported components reachable
  by other applications.
\end{itemize}
Signature-based screening (rule sets expressing known fraud patterns) and a mapping of
findings to a standard adversary-technique taxonomy (MITRE ATT\&CK for Mobile) will be
applied to make the output actionable for analysts.

\subsection{Stage 3 --- Controlled dynamic analysis}
Because sophisticated samples conceal behaviour from static inspection (e.g.\ encrypted
payloads decrypted only at runtime), we will optionally execute samples inside an
\emph{isolated, instrumented sandbox}. The sandbox will install the application, exercise
all of its entry points (launcher activity, broadcast receivers, and background
services), and observe runtime behaviour: granted permissions, started services,
contacted network endpoints, file activity, and security-relevant API invocations.

To counter anti-analysis behaviour, the sandbox will apply \emph{evasion-resistance}
techniques---presenting realistic device characteristics so that emulator-detection
checks do not suppress malicious behaviour---and will capture dynamically loaded code and
command-and-control endpoints. All execution is confined to a disposable, isolated
environment and is strictly opt-in; the static pipeline never executes the sample.

\subsection{Stage 4 --- Feature engineering}
Static and dynamic observations will be transformed into a fixed-length feature
representation for the classifier. The representation will be drawn from features that
are expressed consistently across data sources, including:
\begin{itemize}
  \item permission presence vectors;
  \item suspicious-API presence/frequency vectors;
  \item intent and component indicators;
  \item counts of network indicators and obfuscation signals;
  \item selected runtime-behaviour and network-flow features (when dynamic analysis is
  available).
\end{itemize}
We will favour features that are robust to superficial obfuscation and that remain
comparable between training data and live submissions, in order to minimise feature
drift between offline training and deployment.

\subsection{Stage 5 --- Machine-learning classification}
We will train supervised classifiers to estimate the probability that a sample is
malicious and, where data permits, its likely family/category (e.g.\ banking, SMS,
adware). Candidate models include:
\begin{itemize}
  \item \textbf{Gradient-boosted trees} (e.g.\ XGBoost/LightGBM) and \textbf{Random
  Forests}, chosen for strong tabular-data performance and interpretability through
  feature-importance analysis;
  \item \textbf{Logistic Regression} as a transparent baseline;
  \item Optional advanced directions, such as graph-based models over the application
  call graph, as a future extension.
\end{itemize}
The classifier output is treated as one \emph{corroborating signal} fused with the
rule-based score---not as a stand-alone verdict---so that the final assessment remains
explainable.

\subsection{Stage 6 --- Hybrid risk scoring}
The final risk score will combine the transparent rule-based score, the machine-learning
probability, and threat-intelligence corroboration:
\[
\text{Risk} \;=\; \text{RuleScore} \;+\; \text{ML adjustment} \;+\;
\text{Intelligence corroboration},
\]
normalised to a $0$--$100$ scale and mapped to severity bands (Low, Medium, High,
Critical). Every contribution is itemised so an analyst can see exactly \emph{why} a
sample received its score. Confirmed multi-engine detections raise confidence, while
patterns characteristic of benign advertisement-heavy applications are down-weighted to
control false positives.

\subsection{Stage 7 --- Generative-AI reasoning and reporting}
A generative-AI layer will consume only the \emph{extracted findings} (never the raw
application) and produce: an executive summary for non-technical fraud managers, a
technical narrative correlating the evidence into a likely attack chain, a list of
indicators of compromise, and recommended actions for fraud/SOC teams. The model is
prompted to reason critically about the evidence---distinguishing, for instance, benign
library usage from genuine malicious intent---so that the report adds analytical value
rather than merely restating the findings.

\subsection{Stage 8 --- Output and integration}
The platform will present an analyst dashboard (risk gauge, itemised breakdown, evidence
panels, and the AI report) and will export machine-readable artefacts: a formal PDF
investigation report and IOC lists suitable for ingestion into blocklists and SOC
tooling.

% ===========================================================================
\section{Evaluation Methodology}
% ===========================================================================
We will report standard classification metrics---precision, recall, F1-score, and
ROC-AUC---computed on held-out data, with particular attention to:
\begin{itemize}
  \item \textbf{False-positive rate on modern benign apps}, since over-flagging
  legitimate applications erodes analyst trust;
  \item \textbf{Temporal (drift) evaluation}, training on earlier samples and testing on
  later ones;
  \item \textbf{Cross-dataset generalisation}, training on one corpus and evaluating on
  another;
  \item \textbf{Explainability checks}, verifying that high-scoring decisions are
  supported by human-interpretable evidence.
\end{itemize}

% ===========================================================================
\section{Ethical and Safety Considerations}
% ===========================================================================
All potentially malicious samples will be handled and executed only within isolated
environments using no personal accounts or networks. The system is designed for
\emph{defensive} purposes---first-level triage and investigation support for banking
fraud and SOC teams. Reporting will be faithful to the evidence, and limitations
(e.g.\ concept drift, sandbox evasion, payloads gated behind unavailable infrastructure)
will be stated transparently rather than overstated.

\end{document}
